% =====================================================================
% 3. Training loop (T8, OF + MM) and evaluation module (T7, MM)
% =====================================================================
\section{Training and evaluation code}\label{sec:training}

\subsection{A task-agnostic training loop}

The loop (\texttt{src/train.py}, written by both of us) serves every classifier and the generators,
because everything specific to a task comes from a \emph{task} object: the loss, the metrics
(accumulated batch by batch, so that a generator never keeps all its logits in memory) and the
selection metric with its direction (macro AUROC, higher is better; or cross-entropy, lower is
better). Listing~\ref{lst:loop} shows the five steps of an epoch and the evaluation pass.

\lstinputlisting[firstnumber=101,caption={The training epoch and the evaluation pass
(\texttt{src/train.py}, lines 101--139).},label={lst:loop}]{common/code/train_one_epoch.py}

PyTorch details behind these lines:
\begin{itemize}[leftmargin=1.2em,itemsep=1pt]
  \item \texttt{optimizer.zero\_grad()} is needed because PyTorch \emph{adds} new gradients to the
        stored ones; \texttt{loss.backward()} fills \texttt{.grad} of every trainable parameter and
        \texttt{optimizer.step()} uses them.
  \item \texttt{model.train()} and \texttt{model.eval()} switch dropout and batch norm between their two
        behaviours; \texttt{@torch.no\_grad()} on the evaluation pass avoids building the autograd graph
        (less memory, faster).
  \item \texttt{loss.item()} turns the loss into a Python number, so the graph of the batch is not kept
        alive by the running sum. The evaluation sums the loss and divides once at the end, which gives
        the exact mean even if the last batch is smaller.
  \item The optimiser receives only parameters with \texttt{requires\_grad=True}; a frozen backbone
        (Section~\ref{sec:strategies-code}) therefore needs no special code. The run log records total and
        trainable parameters.
  \item Gradient clipping (\texttt{clip\_grad\_norm\_}) rescales the whole gradient when its norm exceeds
        a threshold; it is off for the Dense and CNN blocks and on for the others.
\end{itemize}

The classification loss is a single call (Listing~\ref{lst:bce}); the model outputs logits and the
sigmoid lives inside the loss.

\lstinputlisting[firstnumber=91,caption={Multi-label loss (\texttt{src/tasks.py}, lines 91--95). Class
weights are available but were not used.},label={lst:bce}]{common/code/multilabel_loss.py}

\subsection{Evaluation module (Mar\'ia's block, summary)}

\texttt{src/metrics.py} (task T7, done by Mar\'ia and reviewed by me) computes the AUROC per class with
scikit-learn, the macro AUROC as their plain mean, and per-class F1 thresholds chosen on validation on a
grid from 0.01 to 0.99. Intervals use the patient bootstrap of Listing~\ref{lst:boot}: each resample
draws patients with replacement and repeats every record of a drawn patient as many times as the
patient was drawn. Because the resamples depend only on the patients and the seed, two models
evaluated on the same split get the same resamples, which is what makes the paired comparison possible:
the difference of the two metrics is computed on each resample, and its percentiles give the interval
of the difference. Resamples in which a class has no positive or no negative record are discarded and
counted.

\lstinputlisting[firstnumber=101,caption={Patient bootstrap (\texttt{src/metrics.py}, lines
101--119).},label={lst:boot}]{common/code/patient_bootstrap.py}

\subsection{Run logs and numbers}

Every run writes \texttt{results/logs/<run\_id>.json} with its configuration, data description, model
size, per-epoch curves, best epoch, results per split, paths of the checkpoint and the saved predictions,
and \texttt{meta} (date, commit, whether the code under \texttt{src/} had uncommitted changes,
package versions, device). Each block summarises its runs in one summary log.
\texttt{report/runs.json} maps a short alias to each summary, and \texttt{report/make\_numbers.py}
turns every scalar of those logs into a \LaTeX{} macro (\texttt{\textbackslash cnnFinalTestMacroAurocMean}
expands to \cnnFinalTestMacroAurocMean{}). Only the files listed in \texttt{runs.json} are read, so an
old or smoke run cannot slip into a report, and a name collision stops the script. All the tables of the
blocks are generated from the summaries and contain only macros. These reports therefore contain no
measured number typed by hand; regenerating the logs and \texttt{numbers.tex} updates every report.

The flag for uncommitted code had a bug that the assistant first explained wrongly
(Appendix~\ref{app:ai-code}, Sheet C4): each run wrote its own predictions file before checking
\texttt{git status}, so every run looked ``dirty''. The check now looks only at the code under
\texttt{src/}.
